% Imporditud: tools/import_eksperimendid.py
% Allikas: Eesti füüsikaolümpiaadi eksperimendi ülesannete
%   kogu (Rael Kalda, 2023), ülesanne Ü37, lahendus L37.
\setAuthor{EFO žürii}
\setRound{piirkonnavoor}
\setYear{1998}
\setNumber{G E2}
\setDifficulty{2}
\setTopic{vedelike mehaanika}
\setVahendid{klaasnõu vedelikuga, dünamomeeter, koormis, mõõtejoonlaud}
\setPunktid{10}
\setAllikas{Eksperimendi kogu Ü37, lahendus lk 50}
\setLahendusseis{toores}

\prob{Vedeliku tihedus}
Määrake vedeliku tihedus. Hinnake mõõtmistulemuse täpsust.

\hint

\solu
Teooria: Kehale mõjuvat raskusjõudu F saab mõõta dünamomeetriga. Vedelikus mõjub kehale üleslükkejõud, mille valem on Fu = $\rho$gV , kus $\rho$ on vedeliku tihedus, V keha poolt välja tõrjutud vedeliku ruumala, g vaba langemise kiirendus (1 p). Vedelikus asetsev keha mõjub dünamomeetrile jõuga $F_1$ = F $-$ Fu = F $- \rho$gV (1 p). Vedelikku asetatud keha ruumala on võrdne ruumalaga, mis on keha poolt väljatõrjutud. Keha ruumala on võrdne V = Sh, kus S on anumas oleva vedeliku pinna pindala, h on vedeliku nivoo kõrguse muutus anumas sinna asetatud keha mõjul (1 p). Suhtelise vea valem: E(x) = $\Delta$x/x. Summa (vahe) arvutamise viga: E ($x_1$ $\pm$ $x_2$) = ($\Delta$$x_1$ + $\Delta$$x_2$) / ($x_1$ $\pm$ $x_2$) (1 p). Korrutise (jagatise) arvutamise viga: E ($x_1$/$x_2$) = E ($x_1$ $\cdot$ $x_2$) = E ($x_1$)+E ($x_2$) (1 p).

Mõõtmised: Mõõdame dünamomeetriga kehale mõjuva raskusjõu ja vedelikku asetatud kehale mõjuva jõu (1 p). Mõõdame keha ruumala mõõtes anuma ristlõikepindala ja vedeliku nivoo muutuse (1 p).

Arvutused: Arvutame keha ruumala (1 p). Arvutame valemist $F_1$ = F $- \rho$gV vedeliku tiheduse $\rho$ = (F $-$ $F_1$) /gV (1 p).

Mõõtmisvea hindamine: Vedeliku tiheduse suhteline viga E($\rho$) = E (F $-$ $F_1$) + E(g) + E(S) + E(h). Antud ülesande puhul võime arvestamata jätta g ja S vea, mis on teistest olulisemalt väiksemad. Jõu vea hindamisel võiks võtta $\Delta$F võrdseks 0,5 dünamomeetri vähima skaalajaotisega ja $\Delta$h võrdseks joonlaua vähima jaotisega (1 mm), sest siin segab täpsemat fikseerimist menisk ja erinevast klaasi märgamisest tingitud ebatäpsus. (1 p).
\probend
